\Article{37}
In arts.\ 15.\ 16 we explained a method for determining the precision of observations approximately\footnote{The investigation of the same subject which we had presented in an earlier paper (\textit{Bestimmung der Genauigkeit der Beobachtungen. Zeitschrift für Astronomie und verwandte Wissenschaften} Vol.\ I, p.\ 185) rested on the same hypothesis concerning the nature of the function expressing the probability of errors on which we had based the method of least squares in Theoria motus corporum coelestium (see art.\ 9, III).}. But this method assumes that a sufficiently large number of the errors that actually occurred are known exactly, a condition which, strictly speaking, will very rarely, not to say never, be met. If, indeed, the quantities whose approximate values have become known through observations depend, according to a known law, on one or more unknown quantities, the most plausible values of these unknowns can be found by the method of least squares. The values of the quantities that were the objects of observation, computed from those values, will then be supposed to differ very little from the
\SourceLeaf{052}%
true values, so that the greater the number of observations, the more justification there will be for taking their differences from the observed values as the true errors of the observations. This practice has been followed by all calculators who have undertaken to estimate the precision of observations a posteriori in particular cases: but it is clearly wrong in theory, and although in many cases it may suffice for practical purposes, in others it can lead to enormous errors. This subject therefore very much deserves a more careful examination.

In this investigation we shall retain the notation used from art.\ 19 onward. The practice we have described regards the quantities \(A,B,C\) etc. as the true values of \(x,y,z\), and consequently \(\lambda,\lambda',\lambda''\) etc. as the true values of the functions \(v,v',v''\) etc. If all observations have equal precision and their weight \(p=p'=p''\) etc. has been taken as unity, the same quantities, with their signs changed, represent the errors of the observations under that assumption; the rules of art.\ 15 then give the mean error of the observations \(m\)
\[
=\BareRoot{\frac{\lambda\lambda+\lambda'\lambda'+\lambda''\lambda''+\text{etc.}}{\pi}}=\BareRoot{\frac{M}{\pi}}
\]
If the observations do not have the same precision, the quantities \(-\lambda,-\lambda',-\lambda''\) etc. represent the errors of the observations multiplied by the square roots of their weights, and the rules of art.\ 16 would lead to the same formula \(\BareRoot{\frac{M}{\pi}}\), now denoting the mean error of observations assigned weight \(=1\). But clearly an exact calculation would require that, instead of the quantities \(\lambda,\lambda',\lambda''\) etc., the values of the functions \(v,v',v''\) etc. obtained from the true values of \(x,y,z\) etc. be used; i.\ e. instead of \(M\), the value of the function \(\Omega\) corresponding to the true values of \(x,y,z\) etc. Although this value cannot be specified, we are nevertheless certain that it is greater than \(M\) (which is the smallest possible value), except in the infinitely improbable case where the most plausible values of the unknowns agree exactly with the true values. In general, therefore, we can assert that the usual practice produces a mean error smaller than it should be, or assigns the observations too great a precision. Let us now see what rigorous theory teaches.

\Article{38}
First of all it is necessary to investigate how \(M\) depends on the true errors of the observations. Let us denote these, as in art.\ 28, by \(e,e',e''\) etc., and for greater simplicity put
\SourceLeaf{053}
\[
\begin{gathered}
e\BareRoot{p}=\varepsilon,\quad e'\BareRoot{p'}=\varepsilon',\quad e''\BareRoot{p''}=\varepsilon''\quad\text{etc., and also}\\
m\BareRoot{p}=m'\BareRoot{p'}=m''\BareRoot{p''}\quad\text{etc.}=\mu
\end{gathered}
\]
Further, let the true values of \(x,y,z\) etc. be, respectively, \(A-x^0,\allowbreak B-y^0,\allowbreak C-z^0\) etc., and let the corresponding values of \(\xi,\eta,\zeta\) etc. be \(-\xi^0,-\eta^0,-\zeta^0\) etc. Clearly the same values will give for \(v,v',v''\) etc. the values \(-\varepsilon,-\varepsilon',-\varepsilon''\) etc., so that the following holds
\[
\begin{aligned}
\xi^0&=a\varepsilon+a'\varepsilon'+a''\varepsilon''+\text{etc.}\\
\eta^0&=b\varepsilon+b'\varepsilon'+b''\varepsilon''+\text{etc.}\\
\zeta^0&=c\varepsilon+c'\varepsilon'+c''\varepsilon''+\text{etc.}
\end{aligned}
\]
etc., and also
\[
\begin{aligned}
x^0&=\alpha\varepsilon+\alpha'\varepsilon'+\alpha''\varepsilon''+\text{etc.}\\
y^0&=\beta\varepsilon+\beta'\varepsilon'+\beta''\varepsilon''+\text{etc.}\\
z^0&=\gamma\varepsilon+\gamma'\varepsilon'+\gamma''\varepsilon''+\text{etc.}
\end{aligned}
\]
Finally, let us put
\[
\Omega^0=\varepsilon\varepsilon+\varepsilon'\varepsilon'+\varepsilon''\varepsilon''+\text{etc.}
\]
so that \(\Omega^0\) is equal to the value of the function \(\Omega\) corresponding to the true values of \(x,y,z\) etc. Since the following holds indefinitely
\[
\Omega=M+(x-A)\xi+(y-B)\eta+(z-C)\zeta+\text{etc.}
\]
it also follows that
\[
M=\Omega^0-x^0\xi^0-y^0\eta^0-z^0\zeta^0-\text{etc.}
\]
It is therefore clear that \(M\), when expanded, is a homogeneous function of the second order of the errors \(e,e',e''\) etc., which may become larger or smaller according to the different values of the errors. But insofar as the sizes of the errors remain unknown to us, it will be appropriate to consider this function indefinitely and, above all, to assign its mean value according to the principles of the calculus of probability. We shall find this by writing, in place of the squares \(ee,e'e',e''e''\) etc., respectively, \(mm,m'm',m''m''\) etc., while omitting the products \(ee',ee'',e'e''\) etc. altogether; or, what is the same, by writing in place of each square \(\varepsilon\varepsilon,\varepsilon'\varepsilon',\varepsilon''\varepsilon''\) etc. the value \(\mu\mu\), and disregarding the products \(\varepsilon\varepsilon',\varepsilon\varepsilon'',\varepsilon'\varepsilon''\) etc. entirely. In this way the term \(\Omega^0\) clearly gives \(\pi\mu\mu\); the term \(-x^0\xi^0\) will give
\[
-(a\alpha+a'\alpha'+a''\alpha''+\text{etc.})\mu\mu=-\mu\mu
\]
\SourceLeaf{054}%
and similarly each of the remaining parts will give \(-\mu\mu\), so that the total mean value becomes \(=(\pi-\rho)\mu\mu\), where \(\pi\) denotes the number of observations and \(\rho\) the number of unknowns. The true value of \(M\) may indeed be greater or smaller than the mean, according to the errors chance has produced; but the greater the number of observations, the less significant the discrepancy will be, so that the following can be taken as an approximate value of \(\mu\):
\[
\BareRoot{\frac{M}{\pi-\rho}}
\]
Thus the value of \(\mu\) resulting from the erroneous practice discussed in the preceding article must be increased in the ratio of the quantity \(\BareRoot{(\pi-\rho)}\) to \(\BareRoot{\pi}\).

\clearpage
